% =============================================================================
%  Chapter 16 — Adaptive Nodal Projection for Material Fields
% =============================================================================

\chapter{Adaptive Nodal Projection for Material Fields}
\label{ch:adaptive_projection}

This chapter refines the earlier diagnosis of the TET10 post-processing
artefact and documents the resulting architectural change in
\texttt{VTKModelExporter}.  The key conclusion is that \emph{negative
quadrature weights are not the only failure mode}.  For quadratic simplex
elements, the row-sum lumped nodal measures can be non-positive even under
all-positive quadrature.  Therefore the exporter now chooses the projection
operator adaptively.


% ─────────────────────────────────────────────────────────────────────────
\section{Why the Previous Diagnosis Was Incomplete}
\label{sec:adaptive_prev_diag}

\cref{ch:negative_weights} correctly identified a catastrophic failure mode:
if a simplex rule contains negative weights, then the signed support of the
quadratic simplex basis can drive the lumped $L^2$ denominator to zero or to
the wrong sign.

However, the ParaView artefact observed for the current TET10 path persisted
even after switching to all-positive quadrature rules such as:
\begin{itemize}[nosep]
  \item the 4-point Hammer--Marlowe--Stroud rule
        (\texttt{SimplexIntegrator<3,2>}), and
  \item the 27-point Stroud conical product rule
        (\texttt{ConicalProductIntegrator<3,3>}).
\end{itemize}

That observation immediately rules out ``negative quadrature weights'' as the
sole root cause.  The actual issue lies in the \emph{row-sum lumping itself}.


% ─────────────────────────────────────────────────────────────────────────
\section{Exact Row-Sum Lumping for TET10}
\label{sec:adaptive_tet10_exact}

For the quadratic tetrahedron, the vertex and edge shape functions are:
\begin{align}
  N_i^{(v)} &= \lambda_i(2\lambda_i - 1), \\
  N_{ij}^{(e)} &= 4\lambda_i\lambda_j,
\end{align}
where $\lambda_i$ are barycentric coordinates.

The lumped nodal measure attached to node $I$ is
\begin{equation}
  m_I^{(\ell)} = \int_T N_I \,\dd V.
\end{equation}

Using the standard simplex moment formula
\begin{equation}
  \int_T \lambda_1^{a_1}\lambda_2^{a_2}\lambda_3^{a_3}\lambda_4^{a_4}\,\dd V
  = |T|\,
    \frac{a_1!a_2!a_3!a_4! \, 3!}
         {(a_1+a_2+a_3+a_4+3)!},
\end{equation}
we obtain:

\subsection{Vertex nodes}
\begin{align}
  \int_T N_i^{(v)}\,\dd V
  &= \int_T \left( 2\lambda_i^2 - \lambda_i \right)\,\dd V \\
  &= 2\frac{|T|}{10} - \frac{|T|}{4} \\
  &= -\frac{|T|}{20}.
\end{align}

\subsection{Edge-midpoint nodes}
\begin{align}
  \int_T N_{ij}^{(e)}\,\dd V
  &= 4 \int_T \lambda_i\lambda_j \,\dd V \\
  &= 4\frac{|T|}{20} \\
  &= \frac{|T|}{5}.
\end{align}

Therefore, for \emph{every} non-degenerate TET10:
\begin{equation}
  m_i^{(\ell)} = -\frac{|T|}{20} < 0
  \qquad \text{for the four vertex nodes,}
\end{equation}
and
\begin{equation}
  m_{ij}^{(\ell)} = \frac{|T|}{5} > 0
  \qquad \text{for the six edge nodes.}
\end{equation}

This is an exact statement.  It does not depend on the particular positive
quadrature rule, provided the rule is at least degree-2 exact.


% ─────────────────────────────────────────────────────────────────────────
\section{Consequence: Lumped \texorpdfstring{$L^2$}{L2} Is Not a Nodal Average}
\label{sec:adaptive_not_average}

For a nodal projection to behave as an average, the local weights attached to
each node must be strictly positive.  TET10 violates that requirement at its
vertex nodes by construction.

This has two implications:
\begin{itemize}[nosep]
  \item The projected nodal value is no longer an average of nearby Gauss-point
        data.
  \item The operator behaves like a signed extrapolation, which can overshoot
        even when the Gauss-point field is smooth and monotone.
\end{itemize}

This explains the visual symptom seen in ParaView:
\begin{itemize}[nosep]
  \item the displacement field is correct, because it is a primary nodal field;
  \item the Gauss-point strains are correct, because they come directly from
        the element kinematics;
  \item the \emph{projected} nodal strains are oscillatory for TET10, because
        the transfer operator is not a proper averaging operator.
\end{itemize}

Hence the issue is fundamentally a \emph{post-processing transfer} problem,
not a stiffness integration problem.


% ─────────────────────────────────────────────────────────────────────────
\section{Why HEX27 Remains Safe}
\label{sec:adaptive_hex27}

For the quadratic tensor-product hexahedron, the shape functions are products
of one-dimensional quadratic Lagrange polynomials on $[-1,1]$.  Their lumped
nodal measures remain strictly positive under the standard
$3\times3\times3$ Gauss--Legendre rule used in \texttt{fall\_n}.  Therefore
the row-sum lumped projection still behaves as a genuine nodal average and
produces smooth fields.

This distinction is important:
\begin{center}
\renewcommand{\arraystretch}{1.2}
\begin{tabular}{lcc}
  \toprule
  Family & Local lumped nodal weights & Safe to use lumped $L^2$ by default? \\
  \midrule
  HEX27 & strictly positive & yes \\
  TET10 & mixed sign        & no  \\
  \bottomrule
\end{tabular}
\end{center}


% ─────────────────────────────────────────────────────────────────────────
\section{Adaptive Policy Implemented in \texttt{VTKModelExporter}}
\label{sec:adaptive_policy}

The exporter now exposes:
\begin{lstlisting}
enum class MaterialFieldProjection {
    Auto,
    LumpedL2,
    PolynomialPatchRecovery,
    VolumeWeightedAveraging
};
\end{lstlisting}

The production path is \texttt{Auto}.  Its algorithm is:
\begin{enumerate}[nosep]
  \item For each element geometry, compute the local lumped nodal weights
        \[
          m_i^{(\ell)} \approx \sum_g N_i(\bxi_g)\, w_g\, |J_g|.
        \]
  \item If \emph{all} local weights are strictly positive for \emph{all}
        elements, use lumped $L^2$ projection.
  \item Otherwise, use polynomial patch recovery
        (\texttt{patch\_recover\_to\_nodes()}).
  \item Retain volume-weighted averaging only as an explicit low-order
        fallback path.
\end{enumerate}

This choice is conservative and appropriate for production visualisation:
\begin{itemize}[nosep]
  \item it preserves the smoother lumped $L^2$ path for element families where
        it is mathematically sound;
  \item it automatically protects higher-order simplex elements without forcing
        the user to know the details of the underlying quadrature rule;
  \item it avoids the excessive diffusion of a pure volume-average fallback by
        using a real local reconstruction on unsafe patches.
\end{itemize}

\paragraph{Complexity.}
The detection pass costs
$O(n_{\mathrm{elem}}\,n_{\mathrm{gp}}\,n_{\mathrm{node}})$, which is of the
same order as the projection itself and negligible compared with an FE solve.
Because this happens only during post-processing, the architectural priority is
robustness and traceability, not micro-optimisation.


% ─────────────────────────────────────────────────────────────────────────
\section{Code-Level Traceability}
\label{sec:adaptive_traceability}

The implementation is localised in:
\begin{itemize}[nosep]
  \item \texttt{src/post-processing/VTK/VTKModelExporter.hh}
        — \texttt{MaterialFieldProjection},
        \texttt{lumped\_projection\_weights\_into()},
        \texttt{has\_strictly\_positive\_lumped\_projection\_weights()},
        \texttt{patch\_recover\_to\_nodes()}, and adaptive dispatch in
        \texttt{compute\_material\_fields()}.
  \item \texttt{main.cpp}
        — unchanged call site: \texttt{compute\_material\_fields()},
        but its semantics are now adaptive rather than hard-coded lumped $L^2$.
  \item \texttt{tests/test\_simplex.cpp}
        — regression tests that prove:
        \begin{itemize}[nosep]
          \item TET10 keeps negative vertex lump weights under positive
                quadrature;
          \item HEX27 keeps strictly positive lump weights.
          \item The new patch-recovery path reproduces exact linear and
                quadratic fields on TET10 patches.
        \end{itemize}
\end{itemize}


% ─────────────────────────────────────────────────────────────────────────
\section{Verification}
\label{sec:adaptive_verification}

Two levels of verification were added:
\begin{enumerate}[nosep]
  \item \textbf{Mathematical regression:} the simplex test suite now checks the
        exact TET10 lumped nodal weights under the 4-point positive HMS rule,
        obtaining
        $-\frac{1}{120}$ at the four vertices and $\frac{1}{30}$ at the six
        edge nodes of the reference tetrahedron.
  \item \textbf{Implementation regression:} the same suite verifies that the
        27-point Stroud conical product rule also yields negative TET10 vertex
        lumps despite its strictly positive quadrature weights, while HEX27
        remains strictly positive.  It also verifies exact linear recovery on a
        single TET10 and exact quadratic recovery on a multi-element TET10
        patch.
\end{enumerate}

These checks establish the core architectural point:
\begin{quote}
For TET10, the unsafe behaviour is not caused by the constitutive model, nor by
the computed Gauss-point strains, nor by the positivity of the stiffness
quadrature weights.  It is caused by using row-sum lumped nodal projection on a
shape-function family whose local lumped nodal measures are signed.
\end{quote}
